\documentclass{article}
\usepackage[zh]{homework}

\config{分析\,-\,0}{2026 春季}{4}

\begin{document}

\begin{problem}[2-13]
    构造一个实数的紧集, 使它的极限点构成一个可数集.
\end{problem}

\begin{solution}
    事实上, 可以构造一个实数的紧集, 使它的极限点构成一个单点集. 例如, 集合 $A = \{0\} \cup \{1/n : n \in \mathbb{N}\}$ 就是一个满足条件的集合. 这个集合是紧的, 因为它是闭的且有界的. 它的极限点只有一个, 就是 $0$, 因为 $1/n$ 随着 $n$ 的增大趋近于 $0$. 因此, 这个集合满足题目的要求.
\end{solution}

\begin{problem}[2-15]
    证明: 如果把“紧的”这个词换成“闭的”或“有界的”, 那么定理2.36和它的推论均不成立.
\end{problem}

\begin{note}
    事实上, 应当增加条件: 两个命题中的紧子集组都是无穷的.
\end{note}

\begin{solution}
    抄写一遍定理2.36及其推论:\par
    \textbf{定理2.36} {\kaishu 如果 \( \{K_\alpha\}_{\alpha \in A} \) 是度量空间 \( X \) 的一组紧子集, 并且 \( \{K_\alpha\} \) 中任意有限个集合的交都不是空集, 那么 \( \bigcap_\alpha K_\alpha \) 也不是空集.}\par
    \textbf{定理2.36的推论} {\kaishu 设 \( \{K_n\} \) 是非空紧集的序列, 并且 \( K_n\supset K_{n+1} \) 对任意的正整数 \( n \) 均成立, 那么 \( \bigcap_n K_n \) 也是非空的.}\par
    (1) 首先说明: 把“紧的”换成“闭的”时, 两个命题都不成立. 这是因为, 取 \( X=\R \) 并配备度量 \( d(x,y)=\abs{x-y} \) , 取 \( f: A\to\R \) 并且使得 \( f \) 无上界. 取 \( K_\alpha=[f(\alpha),+\infty),\ K_n=[n,+\infty) \) .\par
    容易验证这两个构造均满足命题条件, 但不满足命题结论.\par
    (2) 把“紧的”换成“有界的”时, 两个命题也都不成立. 这是因为, 取 \( A \) 的一个可数子集 \( A':=\{\alpha_n\}_{n\ge 1} \). 令 \( \{K_n=(0,1/n)\} \) , 并且取
    \[ K_\alpha=\begin{cases}
    (0,1) & \alpha\not\in A', \\
     (0,1/n) & \alpha=\alpha_n\in A'.
    \end{cases} \]  \par
    容易验证这两个构造均满足命题条件, 但不满足命题结论.
\end{solution}

\begin{problem}[2-16]
    把所有有理数组成的集 $\mathbb{Q}$ 看成度量空间, 而以 $d(p,q)=\abs{p-q}$. 令 $E$ 是所有满足
    $2<p^2<3$ 的 $p\in\mathbb{Q}$ 组成的集, 证明 $E$ 是 $\mathbb{Q}$ 中的有界闭集, 但 $E$ 不是
    紧集, $E$ 是否为 $\mathbb{Q}$ 中的开集?
\end{problem}

\begin{solution}
    首先证明 $E$ 是有界的. 由于 $p^2 < 4$, 我们有 $|p| < 2$. 因此, $E \subset (-2, 2)$, 所以 $E$ 是有界的.\par
    再来证明 \( E \) 是闭集. 采用反证法. 设 \( x \) 是 \( E \) 的一个极限点, 而 \( x\not\in E \) . 分类讨论:
    \begin{enumerate}[label=(\alph*)]
        \item 如果 \( 0\le x^2\le 2 \) . 由 \( \sqrt{2}\not\in\Q \) 知道 \( 0\le x^2<2 \) . 取
        \[ r=\frac{2-x^2}{2\abs{x}+2}<2. \] \par
        我们来证明 \( N_r(x)\cap E=\varnothing \) , 这只需要证明 \( (x\pm r)^2\le 2 \) . 这是因为, 
        \[ (x\pm r)^2=x^2+r(\mp2x+r)<x^2+r(2\abs{x}+2)<x^2+(2-x^2)=2. \] 
        \item 如果 \( 3\le x^2 \) . 由 \( \sqrt{3}\not\in\Q \) 知道 \( x^2>3 \) . 取
        \[ r=\frac{x^2-3}{2\abs{x}}>0. \] 
        我们来证明 \( N_r(x)\cap E=\varnothing \) , 这只需要证明 \( (x\pm r)^2\ge 3 \) . 这是因为,
        \[ (x\pm r)^2=x^2\pm 2xr+r^2>x^2-2r\abs{x}=x^2-(x^2-3)=3. \] 
    \end{enumerate}\par
    再来证明 \( E \) 不是紧的. 取 \( a_1=3/2 \) , \( a_{n+1}=(2a_n+3)/(a_n+2) \). 容易验证 \( \lim_{n\to\infty}a_n=\sqrt{3} \) 且 \( \{a_n\} \) 单调递增. 令
    \[ \mathcal{O}_n=(-a_n,a_n). \] 
    则 \( \{\mathcal{O}_n\} \) 是 \( E \) 的一个开覆盖, 但没有有限子覆盖.\par
    最后, \( E \) 是 \( \Q \) 中的开集. 证明方法与证明其为闭集是完全类似的.
\end{solution}

\begin{problem}[2-19]
    (a) 令 $A$ 及 $B$ 是某度量空间 $X$ 中的不交闭集, 证明它们是分离的.

    (b) 证明不交开集也是分离的.

    (c) 固定 $p\in X$, $\delta>0$, 定义 $A$ 为由满足 $d(p,q)<\delta$ 的一切 $q\in X$ 组成的集,
    而 $B$ 为满足 $d(p,q)>\delta$ 的一切 $q\in X$ 组成的集. 证明 $A$ 与 $B$ 是分离的.

    (d) 证明至少含有两个点的连通度量空间, 是不可数的. 提示: 用 (c).
\end{problem}

\begin{proof}
    (a) 由 \( A \) 与 \( B \) 均为闭集, 知 \( A=\overline{A} \) 且 \( B=\overline{B} \) . 故由 \( A\cap B=\varnothing \) 知 \( A\cap\overline{B}=B\cap\overline{A}=\varnothing \) , 因此两者相离.\par
    (b) 设 \( A \) 和 \( B \) 是两个不交的开集. 若两者不是分离的, 不妨设 \( x\in A\cap\overline{B} \) . 由于 \( A \) 是开集, 存在 \( r>0 \) 使得 \( N_r(x)\subset A \) . 由于 \( x\in\overline{B} \) , 存在 \( y\in N_r(x)\cap B \) . 这与 \( A\cap B=\varnothing \) 矛盾.\par
    (c) 不难知道 \( \overline{A}=\{q\in X:d(p,q)\le\delta\} \) 而 \( \overline{B}=\{q\in X:d(p,q)\ge\delta\} \) . 因此 \( A\cap\overline{B}=B\cap\overline{A}=\varnothing \) , 因此两者相离.\par
    (d) 若 \( X \) 是包含至少两个点的连通度量空间, 则任取 \( x\in X \) , 则集合
    \[ D:=\left\{d(x,y):y\in X-\{x\}\right\}\subset\R_+ \] 
    非空. 若 \( X \) 可数, 则 \( D \) 至多可数, 因此存在 \( r\in \R_+-D \) . 则取
    \[ A=\left\{y\in X:d(x,y)<r\right\},\ B=\left\{y\in X:d(x,y)>r\right\}. \] \par
    由 \( r \) 的取法知 \( X=A\cup B \) , 再由(c)知 \( A, B \) 不连通, 这与 \( X \) 连通矛盾.
\end{proof}

\begin{problem}[2-21]
    令 $A$ 及 $B$ 是某个 $\mathbb{R}^k$ 的分离子集, 设 $\bo{a}\in A$, $\bo{b}\in B$, 且对 $t\in\mathbb{R}^1$ 定义
    \[
        p(t)=(1-t)\bo{a}+t\bo{b}.
    \]
    令 $A_0=p^{-1}(A)$, $B_0=p^{-1}(B)$ [于是, $t\in A_0$ 当且仅当 $p(t)\in A$].\par
    (a) 证明 $A_0$ 与 $B_0$ 是 $\mathbb{R}^1$ 中的分离子集.\par
    (b) 证明存在 $t_0\in(0,1)$ 使得 $p(t_0)\not\in A\cup B$.\par
    (c) 证明 $\mathbb{R}^k$ 的凸子集是连通集.
\end{problem}

\begin{proof}
    (a) 若 \( A_0 \) 与 \( B_0 \) 在 \( \R \) 中不分离, 不妨设 \( t_0\in A_0\cap\overline{B_0} \) . 则 \( p(t_0)\in A \) . \par
    我们来证明 \( p(t_0)\in\overline{B} \) , 这样就可以导出矛盾. 对于任意的 \( r>0 \) , 存在 \( t'\in N_{r/\abs{\bo{a}-\bo{b}}}(t_0)\cap B_0 \) , 因此
    \[ f(t')\in f \left(N_{r/\abs{\bo{a}-\bo{b}}}(t_0)\right)\cap f(B_0)\in N_r(f(t_0))\cap B. \] 
    因此 \( p(t_0) \) 是 \( B \) 的一个极限点, 因此 \( p(t_0)\in\overline{B} \) .\par
    (b) 否则, 对于任意的 \( t\in(0,1) \) , 均有 \( p(t)\in A\cup B \) , 即 \( (0,1)=A_0'\cup B_0' \) , 其中 \( A_0'=A_0\cap(0,1) \) ,  \( B_0'=B_0\cap (0,1) \) . 由于 \( (0,1) \) 是连通的, 故 \( A_0' \) 与 \( B_0' \) 连通, 因此 \( A_0,B_0 \) 连通, 与(a)矛盾.\par
    (c) 设 \( E \) 是 \( \R^k \) 的凸集, 而 \( E \) 不连通. 则存在分离的 \( A, B \subset E \) , 使得 \( A\cap B=\varnothing \) 且 \( A\cup B=E \) . 任取 \( \bo{a}\in A \) , \( \bo{b}\in B \) , 则 \( p(t)\in E=A\cup B \) 对任意的 \( t\in(0,1) \) 都成立. 这与(b)矛盾. 
\end{proof}

\begin{problem}[2-25]
    证明紧致度量空间有可数的拓扑基, 因而是可分的.
\end{problem}

\begin{proof}
    设 \( K \) 是一个紧致度量空间, 则 \( K \) 的任意的无限子集均有极限点. 由习题2-24知 \( K \) 是可分的, 设其可数稠密子集为 \( A=\left\{a_n\right\} \) . \par
    我们证明:
    \[ \mathcal{B}:=\left\{N_r(a_n):\ r\in\Q_+,\ a_n\in A\right\} \] 
    是可数的拓扑基. 首先, 由 \( \Q_+ \) 与 \( A \) 可数, 知 \( \mathcal{B} \) 可数. 因此, 只需要证明:  \( K \) 中的任意开集 \( E \) 均为 \( \mathcal{B} \) 中某些集合的并集.\par
    设 \( x\in E \) . 则存在 \( r_x\in\Q_+ \) 使得 \( N_{r_x}(x)\subset E \) . 由于 \( A \) 在 \( K \) 中稠密, 存在 \( a_{n_x}\in A\cap N_{r_x/2}(x) \) . 则 \( x\in N_{r_x/2}(a_{n_x})\subset N_{r_x}(x)\subset E \) . 因此
    \[ E=\bigcup_{x\in E}N_{r_x/2}(a_{n_x}), \]
    其中 \( N_{r_x/2}(a_{n_x})\in\mathcal{B} \) . 因此 \( \mathcal{B} \) 是 \( K \) 的一个拓扑基.
\end{proof}

\begin{problem}[2-26]
    设 \( X \) 是一个度量空间, 满足: \( X \) 内的任意无限子集均有极限点. 求证: \( X \) 是紧的.
\end{problem}

\begin{proof}
    由前面的若干例题, 知道 \( X \) 是可分的, 因而 \( X \) 具有可数的拓扑基 \( \mathcal{B}:=\left\{B_n:n\ge 1\right\} \) . 下面来证明 \( X \) 是紧致度量空间. \par
    采用反证法. 如果 \( X \) 不是紧致度量空间, 则存在 \( X \) 的一个开覆盖 \( \{O_\alpha\} \) , 使得其没有有限的子覆盖. 对于任意的一个下标 \( \alpha \) , 设 \( S_\alpha\subset\N_+ \) 是这样的一个集合, 使得
    \[ O_\alpha=\bigcup_{n\in S_\alpha}B_n. \] \par
    记 \( S=\bigcup_{\alpha}S_\alpha\subset\N_+ \) , 并且将 \( S \) 的元素记为 \( S=\left\{s_1,s_2,\cdots\right\} \) . 则
    \[ X=\bigcup_{\alpha}O_\alpha=\bigcup_{\alpha}\bigcup_{n\in S_\alpha}B_n=\bigcup_{n\in S}B_n. \tag{\( \ast \)}\]  \par
    若 \( S \) 有限, 则这表明 \( \{O_\alpha\} \) 本身就是有限子覆盖, 问题变为平凡的. 下面假设 \( S \) 无限. \par
    由 \( S \) 的定义, 知对于任意的正整数 \( n \) , 存在 \( \alpha_n \) 使得 \( s_n\in S_{\alpha_n} \) . 则 \( B_{s_n}\subset O_{\alpha_n} \) . 对于任意的正整数 \( n \) , 由 \( \{O_\alpha\} \) 没有有限子覆盖, 知
    \[ X\not\subset\bigcup_{k=1}^{n}O_{\alpha_k}\implies\exists\ p_n\in X-\bigcup_{k=1}^{n}O_{\alpha_k}. \]
    记 \( P=\{p_n:n\in\N_+\} \) . 分情况讨论:
    \begin{enumerate}[label=(\alph*)]
        \item  \( P \) 是有限集. 则存在 \( p\in P \) , 使得 \( \left\{n\in\N_+:p_n=p\right\} \) 是无限集. 取正整数 \( m \) 使得 \( p\in B_{s_m} \) (这由(\( \ast \))是可以做到的). 再取正整数 \( m' \) 使得 \( p_{m'}=p \) 且 \( m'>m \) . 则
        \[ p_{m'}=p\in B_{s_m}\subset\bigcup_{k=1}^{m'} B_{s_k}\subset\bigcup_{k=1}^{m'} O_{\alpha_k}. \] 
        这与我们对于点列 \( \{p_n\} \) 的定义矛盾.
        \item  \( P \) 是无限集. 根据条件,  \( P \) 存在极限点, 记为 \( p^* \) . 由(\( \ast \))知可以取正整数 \( m \) 使得 \( p^*\in B_{s_m} \) . 则根据极限点的定义,  \( B_{s_m} \) 包含了 \( P \) 中无穷多个点. 取正整数 \( m'>m \) 使得 \( p_{s_{m'}}\in B_{s_m} \) . 则
        \[ p_{s_{m'}}\in B_{s_m}\subset\bigcup_{k=1}^{m'} B_{s_k}\subset\bigcup_{k=1}^{m'} O_{\alpha_k}. \] 
        这与我们对于点列 \( \{p_n\} \) 的定义矛盾.
    \end{enumerate}\par
    因此, 无论哪种情况均会得到矛盾. 因此 \( X \) 是紧的.
\end{proof}

\begin{problem}[3-1]
    证明: 序列 \( \{s_n\} \) 的收敛性蕴含着 \( \{|s_n|\} \) 的收敛性. 逆命题对吗?
\end{problem}

\begin{solution}
    设 \( \lim_{n\to\infty}s_n=s \). 我们来证明: \( \lim_{n\to\infty}|s_n|=|s| \). \par
    对于任意的 \( \varepsilon>0 \) , 存在 \( N\in \N_+ \), 使得对于任意的 \( n>N \), 均有 \( \abs{s_n-s}<\varepsilon \). 则 \( \abs{\abs{s_n}-\abs{s}}\le\abs{s_n-s}<\varepsilon \) . 因此命题得证.\par
    命题的逆不成立. 取 \( s_n=(-1)^{n} \), 则有 \( \lim_{n\to\infty}\abs{s_n}=1 \), 但是 \( \lim_{n\to\infty}s_n \) 不存在.
\end{solution}

\begin{problem}[3-2]
    计算
    \[ \lim_{n\to\infty}\left(\sqrt{n^2+n}-n\right). \] 
\end{problem}

\begin{solution}
    注意到
    \[ \sqrt{n^2+n}-n=\frac{n^2+n-n^2}{\sqrt{n^2+n}+n}=\frac{1}{\sqrt{1+1/n}+1}\to \frac{1}{\sqrt{1}+1}=\frac{1}{2}. \] 
\end{solution}

\begin{problem}[3-3]
    设 \( s_1=\sqrt{2} \), 且
    \[ s_{n+1}=\sqrt{2+\sqrt{s_n}}\ ,n\ge 1. \] 
    证明 \( \left\{s_n\right\} \) 收敛, 且在 \( n\ge 1 \) 时 \( s_n<2 \) .
\end{problem}

\begin{proof}
    先证明 \( \left\{s_n\right\} \) 单调递增. 显然有 \( s_2=\sqrt{2+\sqrt{2}}>\sqrt{2}=s_1 \). 若 \( s_{n+1}>s_n \), 则有 \( s_{n+2}=\sqrt{2+\sqrt{s_{n+1}}}>\sqrt{2+\sqrt{s_{n}}}=s_{n+1} \), 故由数学归纳法知 \( \left\{s_n\right\} \) 单调递增. \par
    再证明 \( s_n<2 \) 对于所有的正整数 \( n \) 均成立. 仍然使用归纳法. 显然 \( s_1<2 \). 如果 \( s_n<2 \), 则 \( s_{n+1}=\sqrt{2+\sqrt{s_{n}}}<\sqrt{2+\sqrt{2}}<2. \), 故 \( s_n<2,\ \forall\,n\in\N_+ \).\par
    因此, \( \left\{s_n\right\} \) 是单调递增且有上界的数列, 故必然收敛.
\end{proof}

\begin{problem}[3-6]
    研究级数 \( \sum_{n\ge 1}a_n \) 的敛散性, 其中
    \begin{enumerate}[label=(\alph*)]
        \item \( a_n=\sqrt{n+1}-\sqrt{n} \); 
        \item \( a_n=\left(\sqrt{n+1}-\sqrt{n}\right)/n \);
        \item \( a_n=\left(\sqrt[n]{n}-1\right)^n \);
        \item \( a_n=1/\left(1+z^n\right) \), 其中 \( z\in\C \).
    \end{enumerate}
\end{problem}

\begin{solution}
    记部分和序列为 \( S_n=a_1+\cdots+a_n \). \par
    (a) \( s_n=\sqrt{n+1}-1 \), 故发散;\par
    (b) 注意到
    \[ a_n=\frac{\sqrt{n+1}-\sqrt{n}}{n}=\frac{1}{n \left(\sqrt{n}+\sqrt{n+1}\right)}<\frac{1}{2n \sqrt{n}}. \] 
    即 \( a_n=O(n^{-3/2}) \). 结合 \( \sum_{n\ge 1}a_n \) 为正项级数, 不难知道收敛.\par
    (c) 当 \( n \) 充分大时, 有
    \[ n<\left(\frac{\e}{2}\right)^{\sqrt{n}}<\left(1+\frac{1}{\sqrt{n}}\right)^{\sqrt{n}\cdot\sqrt{n}}\implies1+\frac{1}{\sqrt{n}}>\sqrt[n]{n}\implies a_n<n^{-n/2}<n^{-2}. \] 
    故 \( a_n=O(n^{-2}) \). 结合 \( \sum_{n\ge 1}a_n \) 为正项级数, 不难知道收敛.\par
    (d) 当 \( |z|<1 \) 时, \( a_n\to 1 \), 故发散; 当 \( |z|=1 \) 时, 
    \[ \operatorname{Re} a_n=\frac{1}{2}\left(a_n+\overline{a_n}\right)=\frac{1}{2}\frac{2+z^n+z^{-n}}{(1+z^{n})(1+z^{-n})}=\frac{1}{2}\implies \operatorname{Re} S_n=\frac{n}{2}\implies \sum_{n\ge 1}a_n\text{ 发散.} \] 
    当 \( |z|>1 \) 时, 对于充分大的 \( n \), 有
    \[ \abs{a_n}=\frac{1}{\abs{1+z^n}}<\frac{2}{|z|^n}. \] 
    结合 \( |z|>1 \) 知 \( \sum_{n\ge 1}\abs{a_n} \) 收敛, 即 \( \sum_{n\ge 1}a_n \) 绝对收敛, 故收敛.
\end{solution}



\end{document}
